\scanpage{155}
\mspage{V.27}
% archive sheet number: 154

Donc pour avoir la locale $1$-asphéricité pour $\unsure{\mathbb{L}|L}$ en $x$, il suffit de supposer que l'on a, pour tt tel $G$ :
\[
H^i(\deleted{\unsure{X}}V(\unsure{h|k}),G)=0\quad\text{si}\quad i\unsure{\leq|<}1
\tag*{\unsure{(5.5.1)|(5.5.2)}}
\]
On sait \deleted{\unsure{par suite de|par vertu de}} que cette condition est vérifiée \unsure{les|lors} \deleted{\unsure{pour|que}} $\unsure{\deleted{X(h)}|X(h)}$ \deleted{\unsure{est|sont}} \unsure{des|les} intersections complètes de dim $\geq3$ en $x$, \inserted{(et sans doute dès que leur prof géom en $x$ est $\geq3$)}, lorsque $f$ est plat en $x$, cela signifie aussi que $X_0$ est une int. compl. de dimension $\geq2$ en $x$.

Utilisant le th de locale $1$-asphéricité pour $\mathbb{L}$ des morphismes lisses, on trouve ainsi le

\uline{Théorème} 5.6. Supposons que $f:X\longrightarrow S$ soit plat, que l'ens $Z$ des points \inserted{de $X$} en lesquels $f$ est non lisse est de codimension $\geq2$, et que $X_0$ soit une intersection complète locale en tous ses points. Alors $f$ est localement $1$-asphérique pour $\mathbb{L}_{\unsure{\eta|s}}=\mathbb{P}-\{p\}$. En particulier si $f$ est propre, \inserted{alors,} pour $\forall$ groupe fini $G$ \inserted{d'ordre} premier à $p$, l'hom de spécialisation
\[
H^{\unsure{1|i}}(X_0,G)\unsure{-|\longrightarrow}H^1(X_{\bar{\unsure{\eta|y}}},G)
\]
est un isomorphisme.

\uline{Remarque} 5.7. a) Il semble qu'il doive être possible de remplacer l'hypothèse que $X_0$ est une intersection complète par celle que pour tt $x\in Z$, on ait $\operatorname{prof\,géom}_x(X_0)\geq2$; cela résulterait d'une \unsure{solution|réponse} affirmative à SGA 1 $\overline{\underline{\mathrm{XIV}}}$ 6.3
